% Options for packages loaded elsewhere
% Options for packages loaded elsewhere
\PassOptionsToPackage{unicode}{hyperref}
\PassOptionsToPackage{hyphens}{url}
%
\documentclass[
  english,
  russian,
  12pt,
  a4paper,
  DIV=11,
  numbers=noendperiod]{scrreprt}
\usepackage{xcolor}
\usepackage{amsmath,amssymb}
\setcounter{secnumdepth}{5}
\usepackage{iftex}
\ifPDFTeX
  \usepackage[T1]{fontenc}
  \usepackage[utf8]{inputenc}
  \usepackage{textcomp} % provide euro and other symbols
\else % if luatex or xetex
  \usepackage{unicode-math} % this also loads fontspec
  \defaultfontfeatures{Scale=MatchLowercase}
  \defaultfontfeatures[\rmfamily]{Ligatures=TeX,Scale=1}
\fi
\usepackage{lmodern}
\ifPDFTeX\else
  % xetex/luatex font selection
\fi
% Use upquote if available, for straight quotes in verbatim environments
\IfFileExists{upquote.sty}{\usepackage{upquote}}{}
\IfFileExists{microtype.sty}{% use microtype if available
  \usepackage[]{microtype}
  \UseMicrotypeSet[protrusion]{basicmath} % disable protrusion for tt fonts
}{}
\usepackage{setspace}
% Make \paragraph and \subparagraph free-standing
\makeatletter
\ifx\paragraph\undefined\else
  \let\oldparagraph\paragraph
  \renewcommand{\paragraph}{
    \@ifstar
      \xxxParagraphStar
      \xxxParagraphNoStar
  }
  \newcommand{\xxxParagraphStar}[1]{\oldparagraph*{#1}\mbox{}}
  \newcommand{\xxxParagraphNoStar}[1]{\oldparagraph{#1}\mbox{}}
\fi
\ifx\subparagraph\undefined\else
  \let\oldsubparagraph\subparagraph
  \renewcommand{\subparagraph}{
    \@ifstar
      \xxxSubParagraphStar
      \xxxSubParagraphNoStar
  }
  \newcommand{\xxxSubParagraphStar}[1]{\oldsubparagraph*{#1}\mbox{}}
  \newcommand{\xxxSubParagraphNoStar}[1]{\oldsubparagraph{#1}\mbox{}}
\fi
\makeatother


\usepackage{longtable,booktabs,array}
\usepackage{calc} % for calculating minipage widths
% Correct order of tables after \paragraph or \subparagraph
\usepackage{etoolbox}
\makeatletter
\patchcmd\longtable{\par}{\if@noskipsec\mbox{}\fi\par}{}{}
\makeatother
% Allow footnotes in longtable head/foot
\IfFileExists{footnotehyper.sty}{\usepackage{footnotehyper}}{\usepackage{footnote}}
\makesavenoteenv{longtable}
\usepackage{graphicx}
\makeatletter
\newsavebox\pandoc@box
\newcommand*\pandocbounded[1]{% scales image to fit in text height/width
  \sbox\pandoc@box{#1}%
  \Gscale@div\@tempa{\textheight}{\dimexpr\ht\pandoc@box+\dp\pandoc@box\relax}%
  \Gscale@div\@tempb{\linewidth}{\wd\pandoc@box}%
  \ifdim\@tempb\p@<\@tempa\p@\let\@tempa\@tempb\fi% select the smaller of both
  \ifdim\@tempa\p@<\p@\scalebox{\@tempa}{\usebox\pandoc@box}%
  \else\usebox{\pandoc@box}%
  \fi%
}
% Set default figure placement to htbp
\def\fps@figure{htbp}
\makeatother



\ifLuaTeX
\usepackage[bidi=basic,provide=*]{babel}
\else
\usepackage[bidi=default,provide=*]{babel}
\fi
% get rid of language-specific shorthands (see #6817):
\let\LanguageShortHands\languageshorthands
\def\languageshorthands#1{}


\setlength{\emergencystretch}{3em} % prevent overfull lines

\providecommand{\tightlist}{%
  \setlength{\itemsep}{0pt}\setlength{\parskip}{0pt}}



 
\usepackage[style=gost-numeric,backend=biber,langhook=extras,autolang=other*]{biblatex}
\addbibresource{bib/cite.bib}

\usepackage[]{csquotes}

\usepackage{indentfirst}
\usepackage{float}
\floatplacement{figure}{H}
\IfFileExists{plex-otf.sty}{
  %% Full TeXlive
  \usepackage[
    % math,
    RM={Scale=0.94},SS={Scale=0.94},SScon={Scale=0.94},TT={Scale=MatchLowercase,FakeStretch=0.9},
    DefaultFeatures={Ligatures=Common}
  ]{plex-otf}
  % \usepackage[Scale=MatchUppercase]{juliamono}
}{
  %% TinyTeX
  \usepackage{libertine}
}
\KOMAoption{captions}{tableheading}
\makeatletter
\@ifpackageloaded{caption}{}{\usepackage{caption}}
\AtBeginDocument{%
\ifdefined\contentsname
  \renewcommand*\contentsname{Содержание}
\else
  \newcommand\contentsname{Содержание}
\fi
\ifdefined\listfigurename
  \renewcommand*\listfigurename{Список иллюстраций}
\else
  \newcommand\listfigurename{Список иллюстраций}
\fi
\ifdefined\listtablename
  \renewcommand*\listtablename{Список таблиц}
\else
  \newcommand\listtablename{Список таблиц}
\fi
\ifdefined\figurename
  \renewcommand*\figurename{Рисунок}
\else
  \newcommand\figurename{Рисунок}
\fi
\ifdefined\tablename
  \renewcommand*\tablename{Таблица}
\else
  \newcommand\tablename{Таблица}
\fi
}
\@ifpackageloaded{float}{}{\usepackage{float}}
\floatstyle{ruled}
\@ifundefined{c@chapter}{\newfloat{codelisting}{h}{lop}}{\newfloat{codelisting}{h}{lop}[chapter]}
\floatname{codelisting}{Список}
\newcommand*\listoflistings{\listof{codelisting}{Листинги}}
\makeatother
\makeatletter
\makeatother
\makeatletter
\@ifpackageloaded{caption}{}{\usepackage{caption}}
\@ifpackageloaded{subcaption}{}{\usepackage{subcaption}}
\makeatother
\usepackage{bookmark}
\IfFileExists{xurl.sty}{\usepackage{xurl}}{} % add URL line breaks if available
\urlstyle{same}
\hypersetup{
  pdftitle={Индивидуальный проект},
  pdfauthor={Курушин Георгий Романович},
  pdflang={ru-RU},
  hidelinks,
  pdfcreator={LaTeX via pandoc}}


\title{Индивидуальный проект}
\usepackage{etoolbox}
\makeatletter
\providecommand{\subtitle}[1]{% add subtitle to \maketitle
  \apptocmd{\@title}{\par {\large #1 \par}}{}{}
}
\makeatother
\subtitle{2 Этап}
\author{Курушин Георгий Романович}
\date{}
\begin{document}
\maketitle

\renewcommand*\contentsname{Содержание}
{
\setcounter{tocdepth}{1}
\tableofcontents
}
\listoffigures
\listoftables

\setstretch{1.5}
\chapter{Цель
работы}\label{ux446ux435ux43bux44c-ux440ux430ux431ux43eux442ux44b}

Приобретение практических навыков по установке DVWA

\chapter{Задание}\label{ux437ux430ux434ux430ux43dux438ux435}

Установить дистрибутив DVWA в Kali Linux

\chapter{Теоретическое
введение}\label{ux442ux435ux43eux440ux435ux442ux438ux447ux435ux441ux43aux43eux435-ux432ux432ux435ux434ux435ux43dux438ux435}

DVWA - это уязвимое веб-приложение, разработанное на PHP и MYSQL.

Некоторые из уязвимостей веб приложений, который содержит DVWA: -
Брутфорс: Брутфорс HTTP формы страницы входа - используется для
тестирования инструментов по атаке на пароль методом грубой силы и
показывает небезопасность слабых паролей. - Исполнение (внедрение)
команд: Выполнение команд уровня операционной системы. - Межсайтовая
подделка запроса (CSRF): Позволяет «атакующему» изменить пароль
администратора приложений. - Внедрение (инклуд) файлов: Позволяет
«атакующему» присоединить удалённые/локальные файлы в веб приложение. -
SQL внедрение: Позволяет «атакующему» внедрить SQL выражения в HTTP из
поля ввода, DVWA включает слепое и основанное на ошибке SQL внедрение. -
Небезопасная выгрузка файлов: Позволяет «атакующему» выгрузить
вредоносные файлы на веб сервер. - Межсайтовый скриптинг (XSS):
«Атакующий» может внедрить свои скрипты в веб приложение/базу данных.
DVWA включает отражённую и хранимую XSS. - Пасхальные яйца: раскрытие
полных путей, обход аутентификации и некоторые другие.

DVWA имеет четыре уровня безопасности, они меняют уровень безопасности
каждого веб приложения в DVWA: - Невозможный --- этот уровень должен
быть безопасным от всех уязвимостей. Он используется для сравнения
уязвимого исходного кода с безопасным исходным кодом. - Высокий --- это
расширение среднего уровня сложности, со смесью более сложных или
альтернативных плохих практик в попытке обезопасить код. Уязвимости не
позволяют такой простор эксплуатации как на других уровнях. - Средний
--- этот уровень безопасности предназначен главным образом для того,
чтобы дать пользователю пример плохих практик безопасности, где
разработчик попытался сделать приложение безопасным, но потерпел
неудачу. - Низкий --- этот уровень безопасности совершенно уязвим и
совсем не имеет защиты. Его предназначение быть примером среди уязвимых
веб приложений, примером плохих практик программирования и служить
платформой обучения базовым техникам эксплуатации.
\autocite[\textcite{parasram}]{guide}

\chapter{Выполнение лабораторной
работы}\label{ux432ux44bux43fux43eux43bux43dux435ux43dux438ux435-ux43bux430ux431ux43eux440ux430ux442ux43eux440ux43dux43eux439-ux440ux430ux431ux43eux442ux44b}

Перехожу в директорию \texttt{/var/www/html}. Затем клонирую нужный
репозиторий GitHub (\textbf{?@fig-001}).

\begin{figure}

{\centering \includegraphics[width=0.7\linewidth,height=\textheight,keepaspectratio]{image/1.png}

}

\caption{Клонирование репозитория}

\end{figure}%

Чтобы настроить DVWA, нужно перейти в каталог \texttt{/dvwa/config},
затем проверяю содержимое каталога (\textbf{?@fig-002}).

\begin{figure}

{\centering \includegraphics[width=0.7\linewidth,height=\textheight,keepaspectratio]{image/2.png}

}

\caption{Перемещение по директориям}

\end{figure}%

Изменяю файл, используемый для настройки DVWA
\texttt{config.inc.php.dist} на \texttt{config.inc.php}.

Далее открываю файл в текстовом редакторе nano (\textbf{?@fig-003}).
Изменяю данные об имени пользователя и пароле.

\begin{figure}

{\centering \includegraphics[width=0.7\linewidth,height=\textheight,keepaspectratio]{image/3.png}

}

\caption{Редактирование файла в nano}

\end{figure}%

По умолчанию в Kali Linux установлен mysql, поэтому можно его запустить
без предварительного скачивания, далее выполняю проверку, запущен ли
процесс (\textbf{?@fig-004}).

\begin{figure}

{\centering \includegraphics[width=0.7\linewidth,height=\textheight,keepaspectratio]{image/4.png}

}

\caption{Запуск mysql}

\end{figure}%

Авторизируюсь в базе данных от имени пользователя root
(\textbf{?@fig-005}). Появляется командная строка с приглашением
\enquote{MariaDB}. Далее создаем в ней нового пользователя, используя
учетные данные из файла config.inc.php. Теперь нужно пользователю
предоставить привилегии для работы с этой базой данных.

\begin{figure}

{\centering \includegraphics[width=0.7\linewidth,height=\textheight,keepaspectratio]{image/5.png}

}

\caption{Авторизация в базе данных и предоставление привилегий}

\end{figure}%

Необходимо настроить сервер apache2, перехожу в соответствующую
директорию (\textbf{?@fig-006}). В файле \texttt{php.ini} нужно будет
изменить один параметр, поэтому открываю файл в текстовом редакторе.

\begin{figure}

{\centering \includegraphics[width=0.7\linewidth,height=\textheight,keepaspectratio]{image/6.png}

}

\caption{Перемещение между директориями}

\end{figure}%

В файле параметры allow\_url\_fopen и allow\_url\_include должны быть
поставлены как \texttt{On} (\textbf{?@fig-007}).

\begin{figure}

{\centering \includegraphics[width=0.7\linewidth,height=\textheight,keepaspectratio]{image/7.png}

}

\caption{Редактирование файла}

\end{figure}%

Запускаем службу веб-сервера apache и проверяем, запущена ли служба
(\textbf{?@fig-008}).

\begin{figure}

{\centering \includegraphics[width=0.7\linewidth,height=\textheight,keepaspectratio]{image/8.png}

}

\caption{Запуск службы}

\end{figure}%

Мы настроили DVWA, Apache и базу данных, поэтому открываем браузер и
запускаем веб-приложение, введя 127.0.0/DVWA (\textbf{?@fig-009}).

\begin{figure}

{\centering \includegraphics[width=0.7\linewidth,height=\textheight,keepaspectratio]{image/9.png}

}

\caption{Запуск веб-приложения}

\end{figure}%

Прокручиваем страницу вниз и нажимем на кнопку
\texttt{create\textbackslash{}reset\ database} (\textbf{?@fig-010}).

\begin{figure}

{\centering \includegraphics[width=0.7\linewidth,height=\textheight,keepaspectratio]{image/10.png}

}

\caption{\enquote{Создание базы данных}}

\end{figure}%

Авторизуюсь с помощью предложенных по умолчанию данных
(\textbf{?@fig-011}).

\begin{figure}

{\centering \includegraphics[width=0.7\linewidth,height=\textheight,keepaspectratio]{image/11.png}

}

\caption{Авторизация}

\end{figure}%

Оказываюсь на домашней странице веб-приложения, на этом установка
окончена (\textbf{?@fig-012}).

\begin{figure}

{\centering \includegraphics[width=0.7\linewidth,height=\textheight,keepaspectratio]{image/12.png}

}

\caption{Домашняя страница DVWA}

\end{figure}%

\chapter{Выводы}\label{ux432ux44bux432ux43eux434ux44b}

В ходе работы я приобрел практические навыки по установке уязвимого
веб-приложения DVWA.

\chapter*{Список
литературы}\label{ux441ux43fux438ux441ux43eux43a-ux43bux438ux442ux435ux440ux430ux442ux443ux440ux44b}
\addcontentsline{toc}{chapter}{Список литературы}

\printbibliography[heading=none]





\end{document}
